\chapter{Conclusões}\label{cap:conclusoes}

A integração da interpretação sísmica, do mapeamento de superfícies e da análise morfométrica permitiu caracterizar os depósitos e complexos de transporte de massa na porção norte da Bacia de Santos, reconhecendo diferenças em sua organização estratigráfica, geometria e relação com o substrato. Foram individualizadas cinco unidades: três depósitos, denominados MTD 1, MTD 2 e MTD 3, e dois complexos, MTC 1 e MTC 2, constituídos, respectivamente, por três e quatro eventos. Essa organização evidencia a recorrência de processos de transporte de massa na área e mostra que sua expressão preservada resulta da interação entre episódios de mobilização, sedimentação e deformação da cobertura.

O relevo associado ao Canal de Santos e às estruturas salinas constitui um controle importante sobre a distribuição e a acomodação dos materiais. Os MTDs 1 e 2 apresentam áreas-fonte próximas à calha e geometrias de preenchimento de depressões, com maior espessura em seu interior e adelgaçamento em direção às margens. No MTC 1, a ocupação de paleocanais pelos eventos A e B e da calha atual pelo evento C indica o aproveitamento sucessivo de compartimentos canalizados como áreas receptoras. Essas relações demonstram a importância da configuração do substrato para a deposição, embora não permitam estabelecer a erosão pelo canal como mecanismo desencadeador das rupturas.

A deformação relacionada à tectônica salina também influencia a geometria dos depósitos, especialmente no MTC 1-A e no MTC 2. A associação entre espessamentos, minibacias e altos da cobertura sustenta a interpretação de acomodação diferencial e de limitação local da distribuição dos materiais. Entretanto, o deslocamento dos horizontes por falhas, particularmente no MTD 3 e no MTC 1-A, demonstra que parte da configuração observada incorpora modificações posteriores à deposição. Assim, o relevo atual não pode ser integralmente utilizado como representação das condições existentes durante os movimentos.

A análise morfométrica permitiu distinguir diferentes formas de organização dos corpos. O volume expressivo da porção mapeada do MTC 1-A está associado principalmente à sua abrangência lateral, enquanto o MTC 2-D se destaca pelos fortes contrastes de espessura e pelas acumulações localizadas em compartimentos deprimidos. Nos MTDs 1 e 2, predominam percursos relativamente curtos entre as áreas de ruptura e a calha receptora. Os resultados mostram que extensão, espessura e volume devem ser interpretados conjuntamente, considerando a geometria basal e os limites da cobertura sísmica. Para os eventos A e C do MTC 1, as dimensões calculadas correspondem apenas às porções investigadas, pois ambos ultrapassam a área disponível.

As diferenças de configuração sísmica interna revelam graus variados de continuidade e deformação dos refletores, contribuindo para a caracterização dos depósitos. Contudo, essas configurações não estabelecem uma correspondência exclusiva com mecanismos de transporte. A distinção entre deslizamentos relativamente coerentes, massas deformadas e fluxos mais desagregados depende da análise conjunta das estruturas internas, das superfícies delimitadoras e das relações geométricas. Da mesma forma, a integração dos critérios classificatórios ampliou a comparação entre as unidades, preservando a distinção entre número de eventos, relação com a área-fonte, proporções geométricas e confinamento.

As depressões interpretadas como \textit{pockmarks} apresentam relações estratigráficas distintas com os movimentos de massa. Foram reconhecidas possíveis feições anteriores a determinados eventos, posteriormente recobertas, e depressões desenvolvidas na cobertura de depósitos já instalados. Essas relações são compatíveis com diferentes episódios de circulação de fluidos. Entretanto, sua ocorrência e sua resposta de amplitude não demonstram atividade atual, composição dos fluidos ou sobrepressões capazes de desencadear rupturas. A circulação constitui, portanto, uma hipótese relevante para a evolução da área, cuja participação na instabilidade permanece dependente de controles adicionais.

A aplicação do HVSR aos registros do OBS 18 acrescentou uma estimativa complementar da profundidade de um possível contraste de impedância na subsuperfície rasa. A sobreposição dos envelopes obtidos nas duas datas analisadas e sua compatibilidade com a região basal do evento MTC 2-D sustentam a investigação de uma associação entre a resposta espectral e a organização estratigráfica local. O contato entre os eventos D e C representa uma interface plausível, mas sua identificação não é exclusiva. A interpretação permanece condicionada às velocidades adotadas, à estimativa indireta de Vs, às simplificações do modelo unidimensional e à disponibilidade de registros em um único ponto.

As estimativas cronológicas situaram os topos das unidades em intervalos do Pleistoceno, sob a hipótese de aplicação da taxa média de sedimentação adotada. Esses valores fornecem uma referência temporal exploratória, mas a extrapolação de uma taxa obtida em um registro relativamente curto para períodos muito mais longos impede correlações precisas com ciclos climáticos ou eventos específicos. Nos complexos, a idade estimada para o topo também não corresponde à idade de todos os episódios internos. As relações de sobreposição reconhecidas na sísmica permanecem, assim, como o controle mais consistente da cronologia relativa.

O trabalho demonstra que a distribuição dos movimentos de massa está vinculada à compartimentação geomorfológica e estrutural do talude. Os controles sobre o percurso e a acumulação dos materiais são mais bem estabelecidos que os gatilhos imediatos das rupturas. A caracterização realizada fornece uma base para investigar a evolução sedimentar e a instabilidade submarina nesse setor da Bacia de Santos, sem permitir, isoladamente, determinar a estabilidade atual ou prever novos eventos.

Os objetivos propostos foram atendidos no âmbito da cobertura e da resolução dos dados disponíveis. Assim, O trabalho permitiu identificar, mapear e caracterizar morfometricamente os depósitos e complexos, aplicar critérios classificatórios e investigar suas relações com a deformação salina e as feições interpretadas como \textit{pockmarks}. As estimativas cronológicas e de profundidade obtidas por HVSR forneceram resultados complementares, condicionados às premissas e aos modelos adotados.

\section{Recomendações para trabalhos futuros}

Recomenda-se priorizar a construção de um modelo tridimensional de velocidades de ondas compressionais, incorporando maior número de controles de poços e, quando disponíveis, perfis sônicos, dados de check-shot e informações do processamento sísmico. O modelo deverá considerar as variações verticais e laterais das propriedades sedimentares e a compartimentação estrutural. Esse aprimoramento permitiria avaliar a sensibilidade das espessuras, dos soterramentos e dos volumes ao modelo de conversão, além de reduzir a dependência de velocidades representativas únicas por evento.

Para a interpretação do HVSR, é especialmente importante obter controles locais de velocidade das ondas de cisalhamento. Um modelo de Vs mais bem condicionado permitiria avaliar a associação do pico com as interfaces interpretadas sem depender exclusivamente de sua conversão a partir de Vp e de cenários de coeficiente de Poisson. Recomenda-se ampliar a cobertura de OBS, analisar períodos mais longos e testar modelos estratificados de resposta sísmica, considerando a influência das condições oceanográficas, do acoplamento dos instrumentos ao fundo e da resposta instrumental.

A ampliação da cobertura sísmica é necessária para delimitar integralmente os eventos A e C do MTC 1 e investigar as áreas-fonte ainda não reconhecidas. Dados de maior resolução, associados à batimetria multifeixe, também poderiam esclarecer as terminações dos corpos, as pequenas escarpas e as relações entre falhas e superfícies deposicionais. A restauração estrutural de setores selecionados contribuiria para separar os efeitos do relevo anterior à deposição daqueles produzidos pela deformação posterior.

Recomenda-se obter testemunhos e controles cronológicos em posições estratigráficas selecionadas, buscando restringir as idades dos eventos e reconhecer hiatos, erosão e variações nas taxas de acumulação. A integração com dados sedimentológicos e geotécnicos permitiria investigar contrastes de compactação, níveis de menor resistência, permeabilidade e pressão de poros. Esses parâmetros são necessários para testar hipóteses de ruptura e desenvolver análises de estabilidade compatíveis com as condições locais.

A investigação dos mecanismos desencadeadores deve priorizar a reconstrução das condições anteriores às rupturas. Nos MTDs 1 e 2, recomenda-se procurar evidências de erosão nas margens do Canal de Santos anteriores à instalação dos depósitos, para avaliar a hipótese de remoção de suporte das encostas. No MTD 3 e nos complexos associados ao domínio salino, a análise de estratos de crescimento, das variações de rejeito das falhas e das relações de corte e recobrimento poderá contribuir para estabelecer se a deformação antecedeu, acompanhou ou sucedeu os movimentos.

A integração de testemunhos, perfis geofísicos e ensaios geotécnicos permitiria investigar a existência de camadas mecanicamente frágeis junto às superfícies basais. Recomenda-se determinar a resistência ao cisalhamento, a permeabilidade, a compressibilidade e o estado de adensamento dos sedimentos, além de avaliar sua resposta ao carregamento cíclico. Esses dados possibilitariam testar se as rupturas poderiam ocorrer por aumento de pressão de poros, sobrecarga sedimentar ou perturbações sísmicas, distinguindo condições predisponentes dos processos capazes de iniciar o movimento.

Para avaliar a participação de fluidos, recomenda-se integrar o mapeamento de possíveis condutos e feições de escape com análises de gás e de água intersticial, medidas de pressão de poros e modelagem do fluxo na cobertura sedimentar. A cronologia dessas feições em relação às superfícies de ruptura seria essencial para investigar se a circulação antecedeu os eventos. A hipótese de dissociação de hidratos exigiria, adicionalmente, evidências de sua presença e a reconstrução das condições de pressão e temperatura relevantes.

O aprimoramento das idades dos depósitos também permitiria testar possíveis associações com episódios de sedimentação acelerada, mudanças na circulação oceânica e abalos sísmicos. A identificação de eventos aproximadamente contemporâneos em setores distintos poderia orientar a investigação de uma perturbação regional, embora a contemporaneidade, isoladamente, não comprove um gatilho comum.

Por fim, sugere-se ampliar o mapeamento de \textit{pockmarks} para além dos depósitos e incorporar às análises de incerteza a delimitação dos contornos e as correlações entre superfícies e erros espaciais. Esse conjunto de avanços permitiria avaliar com maior consistência as relações entre transporte de massa, deformação salina e circulação de fluidos, bem como a relevância dessas relações para a suscetibilidade do talude.

\chapter{Agradecimentos}\label{cap:agradecimentos}

Agradeço a ANP (Agência Nacional de Petróleo, Gás Natural e Biocomustíveis) pela disponibilização dos dados sísmicos e de poços utilizados neste trabalho, concedida por meio do processo SEI nº 48610.232732/2025-71, de 03 de dezembro de 2025. Ao projeto SISMO-OCEANO — Monitoramento Sismológico e Oceanográfico de um Segmento na Margem Sudeste do Brasil: Norte da Bacia de Santos ao Sul da Bacia do Espírito Santo, desenvolvido pela Universidade Federal de Santa Catarina (UFSC) em parceria com o Observatório Nacional (ON), com financiamento da Petrobras, agradeço pela disponibilização dos registros dos sismógrafos de fundo oceânico (OBS) utilizados nas análises de HVSR. Agradeço também ao Observatório Nacional e ao Centro de Pesquisas, Desenvolvimento e Inovação da Petrobras (CENPES/Petrobras) pelo acesso ao software \textit{Petrel®}, fundamental para a interpretação dos dados sísmicos e o desenvolvimento desta pesquisa.
